\documentclass[11pt,a4paper]{article}
\usepackage[utf8]{inputenc}
\usepackage[T1]{fontenc}
\usepackage[spanish]{babel}
\usepackage{hyperref}
\usepackage{geometry}
\geometry{margin=2.5cm}
\usepackage{parskip}

\title{008 el alma del objeto}
\author{Alejandro Valenzuela Rivera\\\small Universidad Autónoma de Nuevo León}
\date{\today}
\begin{document}
\maketitle

\begin{quote}
\textbf{Sesgo declarado.} Este trabajo se sostiene desde el sesgo declarado de su autor, y no finge neutralidad. Lo que afirma se puede tumbar, y ese es el punto: el falsificador está escrito, no escondido.
\end{quote}

\section*{Resumen}
El objeto habla desde su propio operador; no se inventa voz.

\section{Por qué se sostiene}
Un objeto con voz inventada es una pantalla.

\section{De dónde viene}
La autoridad, y el repo: los objetos que hablan desde su operador.

\section{El caso}
Ninguna: es la regla de la voz del objeto.

\section{Método}
Un objeto que se puede correr sin su autor, y que sirve igual.

\section{Límites, y qué lo tumba}
Un objeto que solo funciona con quien lo creo presente.

\section{Lo que queda abierto}
\textquestiondown{}Cuándo la voz de un objeto deja de ser suya?

\section*{Cómo se leyó esta tesis}
Esto se escribió en el repositorio man-made-empathy, del que este candado es una pieza: candados/04 el lenguaje/008 el alma del objeto.md. La máquina que lo compuso lo hizo desde este mismo archivo, y no agregó nada que no estuviera escrito. Si algo de esto es falso, se dice en el mismo sitio, y se corrige con fecha.

\end{document}
